\section{Week 6: Evaluating Models I: Goodness of Fit, Overfitting and
Cross-Validation}

\begin{quote}
\textbf{Chapter 3: Evaluation of cognitive models} · Lecture 1 h · Lab 2
h \textbf{Running data:} the Week 2--5 forgetting data, now with four
candidate models: exponential (\(k=2\)), power (\(k=2\)),
exponential-with-asymptote (\(k=3\)) and two-trace (\(k=4\)).
\end{quote}

\begin{center}\rule{0.5\linewidth}{0.5pt}\end{center}

\subsection{Key terms}

Important terms of this lesson, with their Persian equivalents (repeated
even if they appeared in earlier weeks).

{\def\LTcaptype{none} % do not increment counter
\begin{longtable}[]{@{}
  >{\raggedright\arraybackslash}p{(\linewidth - 4\tabcolsep) * \real{0.3333}}
  >{\raggedright\arraybackslash}p{(\linewidth - 4\tabcolsep) * \real{0.3333}}
  >{\raggedright\arraybackslash}p{(\linewidth - 4\tabcolsep) * \real{0.3333}}@{}}
\toprule\noalign{}
\begin{minipage}[b]{\linewidth}\raggedright
English
\end{minipage} & \begin{minipage}[b]{\linewidth}\raggedright
\RL{\CMfa فارسی}
\end{minipage} & \begin{minipage}[b]{\linewidth}\raggedright
\RL{\CMfa توضیح کوتاه}
\end{minipage} \\
\midrule\noalign{}
\endhead
\bottomrule\noalign{}
\endlastfoot
goodness of fit & \RL{\CMfa نیکویی برازش} & \RL{\CMfa اینکه مدل داده‌ی موجود را چقدر
خوب توصیف می‌کند} \\
fit / fitting & \RL{\CMfa برازش} & \RL{\CMfa تنظیم پارامترها تا مدل با داده جور
شود} \\
overfitting & \RL{\CMfa بیش‌برازش} (overfitting) & \RL{\CMfa برازش نویز نمونه به‌جای
الگوی واقعی} \\
forgetting curve & \RL{\CMfa منحنی فراموشی} & \RL{\CMfa کاهش یادآوری با گذشت
زمان} \\
exponential / power function & \RL{\CMfa تابع نمایی} / \RL{\CMfa توانی} & \RL{\CMfa دو شکل
رایج برای منحنی فراموشی و یادگیری} \\
trace plot & \RL{\CMfa نمودار} trace & \RL{\CMfa مقدار نمونه‌ها در برابر شماره‌ی
تکرار} \\
model flexibility / complexity & \RL{\CMfa انعطاف‌پذیری} / \RL{\CMfa پیچیدگی مدل} &
\RL{\CMfa دامنه‌ی الگوهایی که مدل می‌تواند تولید کند} \\
training / test data & \RL{\CMfa داده‌ی آموزش} / \RL{\CMfa آزمون} & \RL{\CMfa داده‌ی برازش}
/ \RL{\CMfa داده‌ی کنارگذاشته برای سنجش} \\
cross-validation (CV) & \RL{\CMfa اعتبارسنجی متقابل} (cross-validation) &
\RL{\CMfa برازش روی بخشی از داده و آزمون روی بخش کنارگذاشته} \\
log-likelihood / NLL & \RL{\CMfa لگاریتم درست‌نمایی} / \RL{\CMfa منفی آن} &
\RL{\CMfa مقداری که بهینه‌سازها کمینه یا بیشینه می‌کنند} \\
likelihood & \RL{\CMfa درست‌نمایی} & \RL{\CMfa احتمال داده به‌عنوان تابعی از
پارامترها} \\
generalization & \RL{\CMfa تعمیم} & \RL{\CMfa پیش‌بینی خوب داده‌ی جدید} \\
parameter & \RL{\CMfa پارامتر} & \RL{\CMfa کمیت قابل تنظیم مدل که بین افراد یا شرایط
فرق می‌کند} \\
SSE / RMSD & \RL{\CMfa مجموع مربعات خطا} / \RL{\CMfa ریشه‌ی میانگین مربع انحراف} &
\RL{\CMfa معیارهای فاصله‌ی داده و پیش‌بینی} \\
noise & \RL{\CMfa نویز} & \RL{\CMfa تغییرات تصادفی در شواهد یا در داده} \\
response bias & \RL{\CMfa سوگیری پاسخ} & \RL{\CMfa تمایل به یک پاسخ، مستقل از
محرک} \\
criterion (c) & \RL{\CMfa ملاک} & \RL{\CMfa آستانه‌ی تصمیم؛ نشان‌دهنده‌ی سوگیری
پاسخ} \\
bias--variance trade-off & \RL{\CMfa بده‌بستان سوگیری و واریانس} & \RL{\CMfa مدل ساده
سوگیر، مدل منعطف پرواریانس} \\
estimate / estimator & \RL{\CMfa برآورد} / \RL{\CMfa برآوردگر} & \RL{\CMfa مقدار به‌دست‌آمده
برای پارامتر} / \RL{\CMfa روش به دست آوردن آن} \\
data leakage & \RL{\CMfa نشت داده} & \RL{\CMfa رسیدن اطلاعات داده‌ی آزمون به
برازش} \\
AIC / BIC & \RL{\CMfa معیار اطلاعاتی آکائیکه} / \RL{\CMfa بیزی} & \RL{\CMfa برازش به‌علاوه‌ی
جریمه برای تعداد پارامتر} \\
benchmark & benchmark \RL{\CMfa (مجموعه‌آزمون استاندارد)} & \RL{\CMfa مجموعه سؤال
برای مقایسه‌ی مدل‌های هوش مصنوعی} \\
benchmark contamination & \RL{\CMfa آلودگی} benchmark & \RL{\CMfa بودن سؤال‌های آزمون
در داده‌ی آموزش مدل} \\
parameter recovery & \RL{\CMfa بازیابی پارامتر} & \RL{\CMfa آیا برازش، پارامترهای
معلوم شبیه‌سازی را برمی‌گرداند} \\
\end{longtable}
}

\begin{center}\rule{0.5\linewidth}{0.5pt}\end{center}

\subsection{Learning outcomes}

By the end of this week, students will be able to:

\begin{enumerate}
\def\labelenumi{\arabic{enumi}.}
\tightlist
\item
  \textbf{Explain} why a good fit is necessary but not sufficient
  evidence for a model, using the concepts of flexibility and
  overfitting.
\item
  \textbf{Demonstrate} overfitting by simulation: training error falls
  monotonically with model complexity while prediction error does not.
\item
  \textbf{Implement} hold-out, \(K\)-fold and leave-one-out
  cross-validation with an out-of-sample log-likelihood score, keeping
  the training and test data properly separated.
\item
  \textbf{Choose} what to hold out (trials, conditions, participants)
  according to the generalization claim being made.
\item
  \textbf{Interpret} cross-validation results, including their
  uncertainty, to compare cognitive models.
\end{enumerate}

\begin{center}\rule{0.5\linewidth}{0.5pt}\end{center}

\subsection{Motivation \& big picture}

In Week 3, Exercise 4, the three-parameter exponential-with-asymptote
model fitted the Week 2 data \emph{better} than the power model
(\(\mathrm{SSE} \approx 0.020\) vs.~\(0.026\)), although the data were
generated by the power model. In Week 3, Part D3, a four-parameter
two-trace model did just as well. If we chose models by fit, we would
choose the wrong one, and we would keep adding parameters forever.

The problem is not specific to forgetting curves. A flexible enough
model can fit any data, including the noise. Roberts and Pashler
famously asked how persuasive a good fit really is (see Further
reading). Their answer: not very, unless we also know what the model
could \emph{not} have fitted.

The way out is to remember \textbf{why} we fit models. We want to
capture the \emph{process} that generated the data, and a model that has
captured the process should \textbf{predict new data} from that process.
Noise, by definition, does not repeat. So the question changes from
``how well does the model fit these data?'' to ``how well does it
predict data it has not seen?'' Cross-validation answers this question
directly with data we already have.

\paragraph{Where is this used?}

\begin{itemize}
\tightlist
\item
  \textbf{Choosing between theories.} Whenever two models of memory,
  decision or learning are compared, a good fit alone is not convincing;
  reviewers ask for cross-validation or a measure that penalizes
  flexibility.
\item
  \textbf{Prediction in clinical psychology and psychiatry}: models that
  predict relapse, treatment response or diagnosis must be validated on
  \emph{new patients}, not on the patients they were fitted to. Many
  published predictors fail exactly this test.
\item
  \textbf{Machine learning}: train/validation/test splits and k-fold
  cross-validation are the standard way to choose models and
  hyperparameters; overfitting is the central danger.
\item
  \textbf{Evaluating AI systems}: benchmark scores are only meaningful
  if the test items were not in the training data (see the AI \& LLMs
  section).
\end{itemize}

\begin{center}\rule{0.5\linewidth}{0.5pt}\end{center}

\subsection{Lecture notes}

\paragraph{In plain words}

A model with enough knobs can pass through every data point, including
the random noise. That does not make it a good model: the noise will be
different next time, so a model that copied the noise will predict
\emph{new} data badly. This is \textbf{overfitting}. The fix is to judge
models by how well they predict data they have \textbf{not} seen.
\textbf{Cross-validation} does this with the data we already have: hide
part of the data, fit the model on the rest, predict the hidden part,
and repeat until every part has been hidden once. The model with the
best predictions on hidden data wins. The week's main message: a better
fit to the \emph{same} data is not evidence for a better model; better
\emph{prediction} of new data is.

\textbf{Suggested timing (60 min):} measures of fit and their limits 10
· overfitting and flexibility 15 · the logic of generalization 5 ·
cross-validation 20 · what to hold out 10.

\subsubsection{1. Measures of goodness of fit}

{\def\LTcaptype{none} % do not increment counter
\begin{longtable}[]{@{}
  >{\raggedright\arraybackslash}p{(\linewidth - 6\tabcolsep) * \real{0.2500}}
  >{\raggedright\arraybackslash}p{(\linewidth - 6\tabcolsep) * \real{0.2500}}
  >{\raggedright\arraybackslash}p{(\linewidth - 6\tabcolsep) * \real{0.2500}}
  >{\raggedright\arraybackslash}p{(\linewidth - 6\tabcolsep) * \real{0.2500}}@{}}
\toprule\noalign{}
\begin{minipage}[b]{\linewidth}\raggedright
Measure
\end{minipage} & \begin{minipage}[b]{\linewidth}\raggedright
Definition
\end{minipage} & \begin{minipage}[b]{\linewidth}\raggedright
Scale
\end{minipage} & \begin{minipage}[b]{\linewidth}\raggedright
Remarks
\end{minipage} \\
\midrule\noalign{}
\endhead
\bottomrule\noalign{}
\endlastfoot
RMSD & \(\sqrt{\mathrm{SSE}/N}\) & Data units & Easy to read; assumes
equal weighting \\
\(R^2\) & \(1 - \mathrm{SSE}/\sum_i (y_i - \bar y)^2\) & 0--1 (for
linear models) & Popular but treacherous for nonlinear models: can be
negative, rewards flexibility, and is dominated by large
between-condition variance (fitting a monotone decline earns a high
\(R^2\) almost automatically) \\
Log-likelihood \(\ell(\hat{\boldsymbol\theta})\), NLL & Week 4 & Nats &
The natural measure for likelihood-based models \\
Deviance \(G^2\) & \(2(\ell_{\text{sat}} - \ell)\) & \ensuremath{\approx} \(\chi^2\) &
Absolute test of misfit (Week 4) \\
\end{longtable}
}

All of these measure \textbf{descriptive adequacy}: how well the model
reproduces the data it was fitted to. None of them accounts for how
easily the model \emph{could} have produced those data.

\subsubsection{2. Overfitting and flexibility}

Write each observation as \textbf{signal + noise}:
\(y_i = g(x_i) + \varepsilon_i\), where \(g\) is the true process. A
fitted model tries to reproduce \(y_i\), but only the \(g\) part will be
there again in new data.

\begin{itemize}
\tightlist
\item
  A model that is \textbf{too simple} misses part of \(g\): systematic
  misfit, or \textbf{underfitting} (high bias).
\item
  A model that is \textbf{too flexible} bends to follow the particular
  \(\varepsilon_i\) of this sample: \textbf{overfitting} (high
  variance). Its fit to the sample improves; its predictions for new
  data get worse.
\end{itemize}

As complexity increases, training error always goes down, while expected
prediction error first decreases and then increases: a U-shape (lab Part
A).

\textbf{Flexibility is more than the number of parameters.} Two models
with the same \(k\) can differ greatly in flexibility, i.e., in the
range of data patterns they can produce. A parameter that enters as an
exponent or a shift usually buys more flexibility than one that scales
the curve. The two-trace model can mimic exponential, power-like and
asymptotic curves, which is exactly why its good fit is unimpressive.
Information criteria (Week 7) count parameters; cross-validation and
Bayes factors (Week 10) respond to actual flexibility.

\begin{quote}
\textbf{Advanced: the bias--variance decomposition.} For squared error,
the expected prediction error at a point \(x\) (over new datasets and
new noise) decomposes as
\[\mathbb{E}\big[(y_{\text{new}} - \hat f(x))^2\big] = \underbrace{\sigma^2}_{\text{irreducible noise}} + \underbrace{\big(g(x) - \mathbb{E}[\hat f(x)]\big)^2}_{\text{bias}^2} + \underbrace{\mathrm{Var}[\hat f(x)]}_{\text{variance}}.\]
Adding parameters lowers bias but raises variance. The best model for
prediction balances the two, and it is not necessarily the true model
when data are scarce. This decomposition is treated in depth in Azimifar
\& Ahmadi (see Further reading).
\end{quote}

\paragraph{Reading the math}

The bias--variance decomposition (Advanced box) splits the average
squared prediction error into three parts:

\begin{itemize}
\tightlist
\item
  \(\sigma^2\), \textbf{irreducible noise}: even the true process cannot
  predict the random part of new data. No model can go below this.
\item
  \textbf{bias²} \(= (g(x) - \mathbb{E}[\hat f(x)])^2\): how far the
  model's \emph{average} prediction (averaged over many possible
  datasets) is from the truth. Too-simple models have large bias.
\item
  \textbf{variance} \(= \mathrm{Var}[\hat f(x)]\): how much the fitted
  curve jumps around from one dataset to the next. Too-flexible models
  have large variance, because they chase each dataset's noise.
\end{itemize}

Adding parameters moves error from the bias term into the variance term.
The best predictor sits where their sum is smallest, which is why a
slightly wrong but simple model can predict better than the true but
complex one when data are scarce.

\subsubsection{3. Generalization as the criterion}

We prefer the model that best predicts \textbf{new data from the same
process}. Three kinds of ``new'' correspond to three scientific claims:

{\def\LTcaptype{none} % do not increment counter
\begin{longtable}[]{@{}
  >{\raggedright\arraybackslash}p{(\linewidth - 4\tabcolsep) * \real{0.3333}}
  >{\raggedright\arraybackslash}p{(\linewidth - 4\tabcolsep) * \real{0.3333}}
  >{\raggedright\arraybackslash}p{(\linewidth - 4\tabcolsep) * \real{0.3333}}@{}}
\toprule\noalign{}
\begin{minipage}[b]{\linewidth}\raggedright
New data are\ldots{}
\end{minipage} & \begin{minipage}[b]{\linewidth}\raggedright
Claim being tested
\end{minipage} & \begin{minipage}[b]{\linewidth}\raggedright
Hold out\ldots{}
\end{minipage} \\
\midrule\noalign{}
\endhead
\bottomrule\noalign{}
\endlastfoot
New trials of the same participant in the same conditions & The model
captures this person's process & Random trials or items \\
New \textbf{conditions} (e.g., longer delays, new stimuli) & The model
captures the mechanism, so it can extrapolate & Whole conditions (the
\emph{generalization criterion} approach) \\
New \textbf{participants} & The model captures what is common to people
& Whole participants \\
\end{longtable}
}

Always state which kind of generalization your evaluation targets.

\subsubsection{4. Cross-validation (CV)}

\textbf{Hold-out validation.} Split the data into a training set (fit)
and a test set (evaluate). Simple, but wasteful and noisy: the result
depends on one particular split.

\textbf{\(K\)-fold CV.} Split the data into \(K\) roughly equal folds.
For each fold \(k = 1, \dots, K\): 1. Fit the model to all data
\textbf{except} fold \(k\), obtaining
\(\hat{\boldsymbol\theta}_{(-k)}\). 2. Score its predictions for fold
\(k\).

Sum the scores over folds. Every observation is used once for testing,
and the model is never tested on data it was trained on. Typical values
are \(K = 5\) or \(10\). \textbf{Leave-one-out (LOO)} is the special
case where every observation is its own fold.

\textbf{The score.} For likelihood-based models, use the
\textbf{out-of-sample log-likelihood} (log predictive score):

\[
\mathrm{CV\text{-}NLL} = -\sum_{k=1}^{K}\ \sum_{i \in \text{fold } k} \ln p\big(y_i \mid \hat{\boldsymbol\theta}_{(-k)}, \mathcal{M}\big).
\]

Lower is better. For Gaussian models, out-of-sample squared error is an
equivalent choice.

\textbf{Rules for honest CV:} 1. \textbf{Everything} that is learned
from data, including parameter estimates, starting-value choices from
grids and data-driven transformations, must be done \emph{inside} each
training fold. Otherwise information ``leaks'' from the test set. 2.
Split along the right unit (§3). For trial-level CV with binomial
counts, split the \textbf{items} within each condition, so that every
training fold still contains every condition. 3. CV estimates are noisy.
Repeat with several random fold assignments and report the mean and
spread, or the SE of the \emph{difference} between models across folds.
4. With the item-level split, CV compares models fitted to \((K-1)/K\)
of the data, so it slightly penalizes models that need a lot of data.

\begin{quote}
\textbf{Advanced: LOO, AIC and Bayes.} For regular models and large
\(N\), leave-one-out CV with the log score and AIC (Week 7) choose the
same model asymptotically. In Bayesian modeling, an efficient
approximation to LOO (PSIS-LOO) is computed from posterior samples. You
will meet it with ArviZ in Week 13.
\end{quote}

\paragraph{Reading the math}

\(\mathrm{CV\text{-}NLL} = -\sum_{k}\sum_{i \in \text{fold } k} \ln p(y_i \mid \hat{\boldsymbol\theta}_{(-k)}, \mathcal{M})\):

\begin{itemize}
\tightlist
\item
  \(\hat{\boldsymbol\theta}_{(-k)}\): the parameters estimated
  \textbf{without} fold \(k\) (the minus sign in the subscript means
  ``leaving out'').
\item
  \(\ln p(y_i \mid \hat{\boldsymbol\theta}_{(-k)})\): how probable the
  model, trained without them, finds the held-out observations. A model
  that predicted them well gives them high probability.
\item
  The double sum adds this over all folds and all held-out points; the
  minus sign makes it a quantity to \emph{minimize}, like the NLL of
  Week 4. It is the same NLL, but each observation is scored by a model
  that never saw it.
\end{itemize}

\subsubsection{5. Interpreting a model comparison}

\begin{itemize}
\tightlist
\item
  Differences in CV-NLL are in \textbf{nats}. As a rough guide,
  differences smaller than about 1--2 nats, or smaller than twice their
  SE across folds or repetitions, should not be over-interpreted.
\item
  A model that wins CV is the best \textbf{predictor} among those
  compared, for this kind of generalization and this amount of data. It
  is not ``the true model''. With more data, more complex models often
  start to win, because their variance shrinks.
\item
  Combine CV with the qualitative checks you already know: residual
  patterns, parameter plausibility, and the model's interpretability.
\end{itemize}

\paragraph{For an engineer}

\begin{itemize}
\tightlist
\item
  \textbf{Overfitting} is the familiar problem of fitting a high-order
  polynomial or a filter with too many taps to a short, noisy record:
  the training error goes to zero and the prediction error explodes.
\item
  The \textbf{bias--variance trade-off} is the same one you meet when
  choosing a model order (AR order, number of filter taps) or a
  smoothing window: too few degrees of freedom gives bias, too many
  gives variance.
\item
  \textbf{Cross-validation} is the data-driven version of model-order
  selection. AIC/MDL-type criteria (Week 7) are the analytic shortcuts
  you may know from AR model order selection.
\item
  \textbf{Training on the test set} is like tuning a receiver on the
  same noise realization you later measure its BER on: the reported
  performance is optimistic.
\item
  \textbf{Nats} are the natural-log version of bits: a difference of 1
  nat is a likelihood ratio of \(e \approx 2.7\).
\end{itemize}

\paragraph{For the instructor}

\begin{itemize}
\tightlist
\item
  \textbf{Start with the polynomial worked example on the board.} Four
  points, four models. Let students pick the best model by training SSE
  (they pick the cubic), then ask what it predicts at 16 days
  (\(-8.5\)). The laugh is the lesson.
\item
  \textbf{Do leave-one-out by hand for one point} (cover the 8-day point
  with your hand, refit the line by eye, predict). Students then
  understand what the code in the lab does.
\item
  \textbf{Insist on ``what is held out''.} Write the three kinds of
  ``new'' data (Section 3) on the board; most mistakes in student
  projects come from holding out the wrong thing.
\item
  \textbf{Typical questions.}

  \begin{itemize}
  \tightlist
  \item
    \emph{``If the true model overfits less, why can a wrong model win
    CV?''} With little data, a simpler wrong model can have lower
    variance. CV picks the best \emph{predictor}, not the true model.
  \item
    \emph{``How many folds?''} 5 or 10 is standard; LOO for very small
    datasets.
  \item
    \emph{``Is \(R^2 = 0.98\) good?''} For a decreasing curve, almost
    any monotone model gets a high \(R^2\); it says little.
  \end{itemize}
\end{itemize}

\begin{center}\rule{0.5\linewidth}{0.5pt}\end{center}

\subsection{Worked example}

\textbf{Data.} Four recall proportions \(y = (0.62, 0.50, 0.47, 0.27)\)
at \(t = (1, 2, 4, 8)\) days. We compare polynomials of increasing
degree as ``models''.

\subsubsection{Step 1: fit to all the data}

{\def\LTcaptype{none} % do not increment counter
\begin{longtable}[]{@{}llll@{}}
\toprule\noalign{}
Model & \(k\) & Fitted equation & Training SSE \\
\midrule\noalign{}
\endhead
\bottomrule\noalign{}
\endlastfoot
Constant & 1 & \(\hat y = 0.465\) & 0.0633 \\
Linear & 2 & \(\hat y = 0.637 - 0.046\,t\) & 0.0032 \\
Quadratic & 3 & \(\hat y = 0.647 - 0.053\,t + 0.0008\,t^2\) & 0.0031 \\
Cubic & 4 & passes exactly through all 4 points & \textbf{0.0000} \\
\end{longtable}
}

By training SSE, the cubic is ``perfect''. But ask it to predict recall
at \(t = 16\) days: it predicts \(-8.5\), i.e., \emph{minus 850\%
recall}. The linear model predicts \(-0.10\) (also impossible, but less
absurd), and the constant predicts 0.47.

\subsubsection{Step 2: leave-one-out CV}

Leave out each point in turn, refit on the remaining three, and predict
the left-out point.

{\def\LTcaptype{none} % do not increment counter
\begin{longtable}[]{@{}
  >{\raggedright\arraybackslash}p{(\linewidth - 8\tabcolsep) * \real{0.2000}}
  >{\raggedright\arraybackslash}p{(\linewidth - 8\tabcolsep) * \real{0.2000}}
  >{\raggedright\arraybackslash}p{(\linewidth - 8\tabcolsep) * \real{0.2000}}
  >{\raggedright\arraybackslash}p{(\linewidth - 8\tabcolsep) * \real{0.2000}}
  >{\raggedright\arraybackslash}p{(\linewidth - 8\tabcolsep) * \real{0.2000}}@{}}
\toprule\noalign{}
\begin{minipage}[b]{\linewidth}\raggedright
Left-out \(t\)
\end{minipage} & \begin{minipage}[b]{\linewidth}\raggedright
Observed
\end{minipage} & \begin{minipage}[b]{\linewidth}\raggedright
Constant predicts
\end{minipage} & \begin{minipage}[b]{\linewidth}\raggedright
Linear predicts
\end{minipage} & \begin{minipage}[b]{\linewidth}\raggedright
Quadratic predicts
\end{minipage} \\
\midrule\noalign{}
\endhead
\bottomrule\noalign{}
\endlastfoot
1 & 0.62 & 0.413 & 0.560 & 0.498 \\
2 & 0.50 & 0.453 & 0.570 & 0.570 \\
4 & 0.47 & 0.463 & 0.448 & 0.330 \\
8 & 0.27 & 0.530 & 0.275 & \textbf{1.250} \\
\textbf{LOO SSE} & & \textbf{0.1125} & \textbf{0.0090} &
\textbf{1.0000} \\
\end{longtable}
}

(The cubic cannot be cross-validated here: three points do not determine
four parameters.)

\textbf{Interpretation.} The quadratic's training SSE is essentially
equal to the linear model's, but its LOO error is more than 100 times
larger. When it does not see the 8-day point, it bends its curvature to
fit the three short delays and predicts recall \emph{above 1} at 8 days.
The extra parameter fitted noise, not forgetting. The linear model
generalizes best of the three. (None of these polynomials is a sensible
\emph{cognitive} model of forgetting, because none respects the 0--1
range; that is exactly why we use exponential and power models instead.)

\begin{center}\rule{0.5\linewidth}{0.5pt}\end{center}

\subsection{Lab (2 hours)}

\subsubsection{Goals}

\begin{enumerate}
\def\labelenumi{\arabic{enumi}.}
\tightlist
\item
  See overfitting happen: training vs.~test error as polynomial degree
  increases.
\item
  Implement item-level \(K\)-fold CV with the out-of-sample binomial
  log-likelihood, and compare four forgetting models.
\item
  Run a \textbf{generalization-criterion} test: train on short delays,
  predict long delays.
\item
  \textbf{Interpret} which model generalizes, and why the best-fitting
  model is not the best predicting one.
\end{enumerate}

\textbf{Timing:} Part A 25 min · Part B 45 min · Part C 25 min · Part D
25 min.

\subsubsection{Setup: data and four models}

\begin{quote}
\textbf{Code · \RL{\CMfa کد:}} \nolinkurl{code/week06.py}, cell 1 · \RL{\CMfa سلول}
1
\end{quote}

\textbf{Expected output:}

\Needspace{7\baselineskip}
\begin{Verbatim}
exp      k = 2  training NLL = 29.812  params = [0.638 0.046]
pow      k = 2  training NLL = 19.830  params = [0.877 0.421]
expasym  k = 3  training NLL = 18.919  params = [0.861 0.38  0.308]
two      k = 4  training NLL = 18.919  params = [0.265 0.    0.596 0.38 ]
\end{Verbatim}

By training NLL, the more complex models win. Notice that the two-trace
model's best fit has a slow component with rate 0: it has collapsed into
the exponential-with-asymptote model (same NLL), as in Week 3. Now let's
see whether they also \emph{predict} better.

\subsubsection{Part A: overfitting with polynomials (25 min)}

We use polynomials in \(x = \ln(1+t)\) as deliberately ``dumb'' flexible
models. The \textbf{test} data are 500 fresh datasets simulated from the
true power model, with the same design.

\begin{quote}
\textbf{Code · \RL{\CMfa کد:}} \nolinkurl{code/week06.py}, cell 2 · \RL{\CMfa سلول}
2
\end{quote}

\textbf{Expected output:}

\Needspace{12\baselineskip}
\begin{Verbatim}
 degree  k  train_MSE  test_MSE
      0  1    0.03980   0.05072
      1  2    0.00663   0.00787
      2  3    0.00204   0.00634
      3  4    0.00202   0.00640
      4  5    0.00173   0.00661
      5  6    0.00087   0.00750
      6  7    0.00087   0.00750
      7  8    0.00000   0.00833
\end{Verbatim}

\textbf{What to notice:} training error decreases with every added
parameter and reaches 0 at \(k = 8\) (interpolation). Test error is
lowest at an intermediate complexity and then rises: the extra
parameters chase noise. Even the best polynomial predicts worse than the
\emph{true} power model, whose test MSE you can compute as an exercise.

\subsubsection{\texorpdfstring{Part B: \(K\)-fold cross-validation over
items (45 min)}{Part B: K-fold cross-validation over items (45 min)}}

Each interval tested 50 items. We assign each item to one of \(K = 5\)
folds \emph{within} each interval, so that every training set contains
all 8 intervals.

\begin{quote}
\textbf{Code · \RL{\CMfa کد:}} \nolinkurl{code/week06.py}, cell 3 · \RL{\CMfa سلول}
3
\end{quote}

\textbf{Expected output} (about 1--2 minutes):

\Needspace{13\baselineskip}
\begin{Verbatim}
         k  train_NLL  CV_NLL_mean  CV_NLL_sd
exp      2     29.812      256.458      1.058
pow      2     19.830      245.791      0.734
expasym  3     18.919      247.370      1.574
two      4     18.919      247.531      1.172

CV-NLL difference to the power model (mean, SD over repeats):
         exp  pow  expasym    two
mean  10.667  0.0    1.579  1.740
std    1.137  0.0    1.807  1.349
\end{Verbatim}

Note that training NLL (binomial, with the binomial coefficients) and
CV-NLL (Bernoulli, item level) are on different scales. Compare models
\textbf{within} a column, not across columns.

\textbf{Interpret:} 1. Which model predicts held-out items best? Is its
advantage over the runner-up larger than the variability across CV
repetitions? 2. The two-trace model has the best training fit. What
happens to it under CV, and why? 3. Can CV tell the power and
exponential-with-asymptote models apart with these data? What would
help?

\subsubsection{Part C: the generalization criterion: predicting longer
delays (25 min)}

A model that has captured \emph{how} forgetting works should be able to
extrapolate. Train on the five shortest delays (\(t \le 7\) days) and
predict the three longest (14, 21, 30 days), which the models have never
seen.

\begin{quote}
\textbf{Code · \RL{\CMfa کد:}} \nolinkurl{code/week06.py}, cell 4 · \RL{\CMfa سلول}
4
\end{quote}

\textbf{Expected output:}

\Needspace{10\baselineskip}
\begin{Verbatim}
         test_NLL_long_delays  pred_t14  pred_t21  pred_t30
model
exp                    72.550     0.086     0.028     0.007
pow                     9.627     0.228     0.186     0.156
expasym                 6.268     0.260     0.257     0.257
two                    46.072     0.131     0.057     0.020
observed: {14: 0.26, 21: 0.26, 30: 0.28}
\end{Verbatim}

\textbf{What to notice:} here the \textbf{exponential-with-asymptote}
model extrapolates best, \emph{not} the true power model. The three
held-out points happen to be flat (0.26, 0.26, 0.28), and a curve that
levels off predicts them well. The power model keeps declining, as the
true process does, but in this sample the noise at the long delays made
recall look flat. The pure exponential and the two-trace model (whose
short-delay fit gives no constraint on the slow component) extrapolate
badly. Lessons: (1) different generalization criteria can rank models
differently; (2) a single test on three data points is itself noisy, so
a model's win on one test set is not proof; (3) agreement across several
criteria (Part B, Part C, Week 7's AIC/BIC) is more convincing than any
one of them.

\subsubsection{Part D: put it together (25 min)}

Fill in this table in your notebook from Parts B and C, and write a
short paragraph (5--8 sentences) answering: \textbf{Which model would
you recommend as the best account of these forgetting data, and how
confident are you?}

{\def\LTcaptype{none} % do not increment counter
\begin{landscape}\small
\begin{longtable}[]{@{}
  >{\raggedright\arraybackslash}p{(\linewidth - 12\tabcolsep) * \real{0.1429}}
  >{\raggedright\arraybackslash}p{(\linewidth - 12\tabcolsep) * \real{0.1429}}
  >{\raggedright\arraybackslash}p{(\linewidth - 12\tabcolsep) * \real{0.1429}}
  >{\raggedright\arraybackslash}p{(\linewidth - 12\tabcolsep) * \real{0.1429}}
  >{\raggedright\arraybackslash}p{(\linewidth - 12\tabcolsep) * \real{0.1429}}
  >{\raggedright\arraybackslash}p{(\linewidth - 12\tabcolsep) * \real{0.1429}}
  >{\raggedright\arraybackslash}p{(\linewidth - 12\tabcolsep) * \real{0.1429}}@{}}
\toprule\noalign{}
\begin{minipage}[b]{\linewidth}\raggedright
Model
\end{minipage} & \begin{minipage}[b]{\linewidth}\raggedright
\(k\)
\end{minipage} & \begin{minipage}[b]{\linewidth}\raggedright
Training NLL
\end{minipage} & \begin{minipage}[b]{\linewidth}\raggedright
5-fold CV-NLL (mean \ensuremath{\pm} SD)
\end{minipage} & \begin{minipage}[b]{\linewidth}\raggedright
Long-delay test NLL
\end{minipage} & \begin{minipage}[b]{\linewidth}\raggedright
Residual pattern (Week 4)
\end{minipage} & \begin{minipage}[b]{\linewidth}\raggedright
Parameters plausible?
\end{minipage} \\
\midrule\noalign{}
\endhead
\bottomrule\noalign{}
\endlastfoot
exp & & & & & & \\
pow & & & & & & \\
expasym & & & & & & \\
two & & & & & & \\
\end{longtable}
\end{landscape}
}

\subsubsection{Exercises (increasing difficulty)}

\textbf{Exercise 1 \ensuremath{\star}.} In Part A, add the true power model with the true
parameters as a ``model'' (it has no fitted parameters). Compute its
test MSE and compare it with the best polynomial's. What is the
irreducible error?

\textbf{Exercise 2 \ensuremath{\star}.} Rerun Part B with \(K = 2\), \(K = 10\) and
leave-one-item-out (\(K = 50\)). Does the ranking change? Which \(K\)
gives the most stable results across repetitions?

\textbf{Exercise 3 \ensuremath{\star}\ensuremath{\star} (data leakage).} Write a ``wrong'' CV in which the
starting values for each fold are taken from the fit to the \emph{full}
dataset, and another in which the fold's test items are accidentally
included in training. Compare the results with the correct CV. Which
mistake matters more, and why?

\textbf{Exercise 4 \ensuremath{\star}\ensuremath{\star} (more data).} Simulate a new dataset from the true
power model with \(n = 500\) items per interval and rerun Parts B and C.
Does the ranking become clearer? Does the two-trace model ever win?

\textbf{Exercise 5 \ensuremath{\star}\ensuremath{\star}\ensuremath{\star} (Advanced: leave-one-participant-out).} Use the
practice data generator from Week 5, Part D, with 20 participants.
Implement CV in which each fold leaves out \emph{one participant}: fit
an aggregate (group-average) model on 19 participants and score the
prediction of the held-out participant's curve with the Gaussian
log-likelihood. Compare the exponential and power models. Then do the
same with individual-level fits. Explain why the two CV schemes answer
different scientific questions.

\subsubsection{Lab check-off}

Show: (1) the training-vs-test plot from Part A, (2) the CV summary
table from Part B, and (3) your Part D paragraph.

\begin{center}\rule{0.5\linewidth}{0.5pt}\end{center}

\subsection{Connection to AI \& LLMs}

\textbf{Benchmark contamination is data leakage.} LLMs are trained on
huge amounts of internet text. If benchmark questions (or their answers)
were published online, they may be in the training data, so the model is
tested on items it has already seen. This is exactly the leakage problem
of cross-validation: the test set is no longer independent of the
training set, so the score measures memorization as well as
generalization. It is one of the most serious problems in evaluating AI
systems, including claims that LLMs show human-like cognitive abilities
on classic psychology tasks, whose materials are widely published.

\begin{quote}
\textbf{Code · \RL{\CMfa کد:}} \nolinkurl{code/week06.py}, cell 5 · \RL{\CMfa سلول}
5
\end{quote}

\textbf{Expected output:}

\Needspace{6\baselineskip}
\begin{Verbatim}
Reported accuracy (all items): 0.73
Accuracy on leaked items:      0.96  (n = 120)
Accuracy on clean items:       0.63  (n = 280)
\end{Verbatim}

The reported score (0.73) overstates the ability to solve \emph{new}
problems (0.63). The fix is the same as in this week's lab: evaluate on
items that were guaranteed not to be used in fitting, for example newly
written items, or systematic variations of classic tasks.

\textbf{Try it (no API needed).} - \ensuremath{\star} Change the share of leaked items
from 0 to 0.9 and plot the reported accuracy. Can you tell from the
reported accuracy alone how much leaked? - \ensuremath{\star}\ensuremath{\star} In practice we do not know
which items leaked. Suppose a ``fresh'' variant of every item (same
logic, new wording) is available. Simulate fresh variants with
\(P(\text{correct}) = 0.60\) and estimate the leakage effect as the
difference between original and fresh accuracy, with a paired bootstrap
CI (Week 5). - \ensuremath{\star}\ensuremath{\star} Discussion: an LLM solves a famous reasoning puzzle
from a 1970s psychology paper. Why is this weak evidence that it
``reasons like a human''? Which test would be stronger?

\begin{center}\rule{0.5\linewidth}{0.5pt}\end{center}

\subsection{Common mistakes \& misconceptions}

\begin{enumerate}
\def\labelenumi{\arabic{enumi}.}
\tightlist
\item
  \textbf{``The model with the best fit is the best model.''} Fit
  rewards flexibility. Always ask how well the model predicts new data.
\item
  \textbf{``Flexibility = number of parameters.''} Functional form
  matters too. Models with equal \(k\) can differ greatly in what they
  can fit.
\item
  \textbf{Leaking test information into training,} e.g., choosing
  starting values, bounds or transformations from the full dataset, or
  normalizing with full-data statistics.
\item
  \textbf{Splitting along the wrong unit.} If the claim is about new
  participants, holding out random trials overestimates generalization.
\item
  \textbf{Over-reading small CV differences.} CV estimates are noisy.
  Repeat the splits and look at the variability of the differences.
\item
  \textbf{``CV picked the true model.''} CV picks the best
  \emph{predictor} for this amount of data. With little data, a simpler
  (even wrong) model can predict better than the true one.
\item
  \textbf{Using \(R^2\) to compare nonlinear cognitive models.} It is
  not comparable across models and rewards flexibility.
\end{enumerate}

\begin{center}\rule{0.5\linewidth}{0.5pt}\end{center}

\subsection{Check your understanding}

\textbf{1.} Why does training error always decrease (or stay equal) when
we add a parameter to a nested model? \textbf{Answer} The larger model
contains the smaller one as a special case (e.g., by setting the extra
parameter to 0), so its best fit can never be worse. The optimizer can
always reproduce the smaller model's fit and usually improves on it by
fitting noise.

\textbf{2.} Define overfitting in terms of signal and noise.
\textbf{Answer} Data = signal (the systematic process) + noise (random
variation). An overfitting model adapts to the noise of the particular
sample. That improves its fit to these data but harms prediction of new
data, in which the noise is different.

\textbf{3.} Describe \(K\)-fold cross-validation in four steps.
\textbf{Answer} (1) Split the data into \(K\) folds. (2) For each fold,
fit the model on the other \(K-1\) folds. (3) Score the fitted model's
predictions for the held-out fold (e.g., with the log-likelihood). (4)
Sum or average the scores over folds; repeat with new random splits if
possible.

\textbf{4.} In the forgetting experiment, why do we assign items to
folds \emph{within} each retention interval rather than leaving out
whole intervals? \textbf{Answer} Within-interval splitting tests
prediction of new trials under the same conditions, and every training
fold still contains every condition. Leaving out whole intervals tests a
different claim: extrapolation to new conditions (the generalization
criterion, Part C).

\textbf{5.} Model A has CV-NLL 2.1 nats lower than model B on average,
but the difference varies with an SD of 3.0 across CV repetitions. What
do you conclude? \textbf{Answer} The evidence is weak. Strictly, the SD
across repetitions reflects split-to-split variability rather than the
sampling uncertainty of the data, but either way a 2-nat difference with
this much instability should not be over-interpreted. Collect more data
or report the models as roughly equally predictive.

\textbf{6.} Why can a simpler, ``wrong'' model sometimes cross-validate
better than the true, more complex model? \textbf{Answer} With limited
data, the complex model's parameters are estimated with high variance,
and that variance hurts prediction more than the simpler model's bias
does (the bias--variance trade-off). With more data, the true model's
variance shrinks and it eventually wins.

\textbf{7.} A student reports \(R^2 = 0.97\) for a nonlinear forgetting
model and concludes that it is correct. Give two objections.
\textbf{Answer} (i) \(R^2\) is high for almost any monotone decreasing
curve, because most variance in forgetting data is the overall decline.
(ii) A good fit says nothing about generalization or flexibility;
residual patterns, CV and a comparison with alternative models are
needed.

\begin{center}\rule{0.5\linewidth}{0.5pt}\end{center}

\subsection{Assignment}

\subsubsection{A3: Model comparison and recovery (1.0 point of 20;
released Week 6, due at the start of the Week 9 lecture)}

A3 spans this week and next: Tasks 1--3 use this week's material, Tasks
4--6 use Week 7's (AIC/BIC, parameter and model recovery).

\textbf{Data.} A forgetting experiment with 12 participants, generated
by this code:

\begin{quote}
\textbf{Code · \RL{\CMfa کد:}} \nolinkurl{code/week06.py}, cell 6 · \RL{\CMfa سلول}
6
\end{quote}

\textbf{Tasks.}

\begin{enumerate}
\def\labelenumi{\arabic{enumi}.}
\tightlist
\item
  \textbf{Individual fits (1.5 pts).} Fit the exponential, power and
  exponential-with-asymptote models to each participant by multi-start
  MLE. Report a table of parameters and NLLs; flag estimates on bounds.
\item
  \textbf{Absolute fit (1 pt).} Compute \(G^2\) per participant and
  model. How many participants does each model fail to fit
  (\(p < .05\))?
\item
  \textbf{Cross-validation (2 pts).} For each participant, run 5-fold
  item-level CV (repeated 5 times) for the three models. Report which
  model wins for how many participants, and the summed CV-NLL across
  participants with its variability.
\item
  \textbf{Information criteria (1.5 pts, Week 7).} Compute AIC, AICc and
  BIC for every participant and model, and the Akaike weights. Do the
  criteria agree with CV?
\item
  \textbf{Recovery (2.5 pts, Week 7).} (a) Parameter recovery for the
  power model at this design: 200 simulations with parameters drawn from
  plausible ranges; plot true vs.~recovered values, report correlations,
  and comment on any trade-off. (b) Model recovery for the three models
  with AIC and BIC: a 3 × 3 confusion matrix from 100 simulations per
  generating model.
\item
  \textbf{Interpretation (1.5 pts, max. 400 words).} Which model best
  describes forgetting in this study, how strong is the evidence, and
  how does the recovery analysis qualify your conclusion?
\end{enumerate}

\textbf{Rubric (10 points, scaled to 1.0):} the points per task are
given above. Within each task: correctness 60\%, clarity of tables and
figures 20\%, interpretation 20\%. Late submissions lose 10\% per day
(max. 3 days).

\begin{center}\rule{0.5\linewidth}{0.5pt}\end{center}

\subsection{Further reading}

\begin{itemize}
\tightlist
\item
  \textbf{Farrell \& Lewandowsky (2018), chapter 10.} Model comparison:
  why fit is not enough, flexibility, and cross-validation in relation
  to other comparison methods. \emph{(Main reading.)}
\item
  \textbf{Azimifar \& Ahmadi (1396)}, vol.~1 (in Persian): the material
  on generalization, overfitting and the bias--variance trade-off from
  the pattern-recognition perspective {[}verify chapter{]}.
\item
  \emph{Optional classics, outside the course list:} Roberts, S., \&
  Pashler, H. (2000). How persuasive is a good fit? A comment on theory
  testing. \emph{Psychological Review, 107}, 358--367. Busemeyer, J. R.,
  \& Wang, Y.-M. (2000). Model comparisons and model selections based on
  generalization criterion methodology. \emph{Journal of Mathematical
  Psychology, 44}, 171--189.
\end{itemize}

\emph{Preview for Week 7:} cross-validation is general but
computationally expensive. Information criteria (AIC, BIC) estimate
generalization from a single fit. And before trusting any comparison, we
must check whether our method can recover the parameters and the models
that actually generated the data.
